\documentclass[a4paper]{article}
\usepackage[utf8]{inputenc}
\usepackage[T2A]{fontenc}
\usepackage[russian]{babel}
\usepackage{amsmath}
\usepackage{amsfonts} 
\usepackage{parskip}
\usepackage{amssymb}
\usepackage{array}
\usepackage[shortlabels]{enumitem}
\usepackage{enumitem}
\usepackage{cancel}
\usepackage{xcolor}
\usepackage{color, soul}
\usepackage{mathbbol}
\usepackage{tikz}
\usepackage{pgfplots}

\pgfplotsset{compat=1.18}
\makeatletter
\AddEnumerateCounter{\asbuk}{\russian@alph}{щ}
\makeatother

\DeclareMathOperator*\lowlim{\underline{lim}}
\DeclareMathOperator*\uplim{\overline{lim}}
% Top Margin
\setlength{\hoffset}{-0.2in}
\setlength{\oddsidemargin}{0cm}
\setlength{\evensidemargin}{0cm}
% Right Margin
\setlength{\voffset}{-0.2in}
\setlength{\topmargin}{0cm}
\setlength{\marginparsep}{0cm}
\setlength{\marginparwidth}{2cm}
\setlength{\headsep}{0.3cm}
\setlength{\headheight}{0.6cm}
% Bottom Margin
\setlength{\textwidth}{16.93cm}
\setlength{\textheight}{24.8cm} % 29.7cm - 2cm\hoffset - 2cm\unten - 0.3cm\headsep - 0.6cm \headheight
\setlength{\footskip}{0cm}
\setlength{\paperheight}{29.7cm}
\newcommand{\Mod}[1]{\ (\mathrm{mod}\ #1)}

\everymath{\displaystyle}
\setul{0.5ex}{0.3ex}
\setulcolor{red}
\renewcommand{\baselinestretch}{1.15} 
\usepackage{graphicx} % Required for inserting images

\begin{document}

\begin{titlepage}
\thispagestyle{empty}

\begin{center}
Национальный исследовательский университет\\
«Высшая школа экономики»

\vspace{3cm}

{\Large\bfseries Анализ предметной области и постановка цели проекта\par}

\vspace{1.5cm}

{\large\bfseries Кейс 1 - Back\par}

\vspace{0.75cm}

Трэвел-приложение для получения «впечатлений» в поездах: создание маршрутов и покупка готовых экскурсий

\vfill

Проектная группа ПМИ
\end{center}

\vspace{1.5cm}

\begin{flushright}
ФИО: \underline{Круковер Денис Ильич}\\[0.5cm]
Группа: \underline{БПМИИ2412}
\end{flushright}
\end{titlepage}

\begin{description}
    \item[\Large1.] \Large \textbf{Анализ предметной области}
    \item[\large1.1] \large \textbf{Описание предметной области}
    
    \normalsize Разрабатываемый продукт представляет собой цифровое трэвел-приложение, ориентированное на формирование и потребление досуга в формате "впечатлений". Продукт позволяет пользователям:

    \begin{itemize}[label=$-$]
        \item самостоятельно создавать маршруты прогулок, поездок и экскурсий
        \item приобретать готовые маршруты и впечатления, созданные авторами
        \item сопровождать своё перемещение в режиме реального времени
        \item получать рекомендации по досугу в зависимости от контекста поездки или города
    \end{itemize}

    Ключевая особенность продукта - фокус не на классическом туризме, а на эмоциональной ценности досуга: \shorthandoff{"}"что делать", "куда пойти", "как провести время с кайфом", в том числе во время пездок на поезде и пребывания в незнакомом городе.

    \item[\large1.2] \large \textbf{Портрет пользователя (персоны)}
    
    \normalsize \textbf{Основная персона}:

    \begin{itemize}[label=$-$]
        \item Возраст: 18-35 лет
        \item Город: крупные города, туристические центры
        \item Контекст использования: поездки, командировки, выходные, спонтанные путешествия
        \item Поведение: \begin{itemize}[label=$\circ$]
            \item Не желание тратить много времени на планирование досуга
            \item Усталость от стандартных туристических маршрутов
            \item Ценит авторский подход и уникальные впечатления
        \end{itemize}
    \end{itemize}

    \textbf{Проблема пользователя}: пользователь не знает, как структурировать досуг, когда есть свободное время, но нет чёткого плана или идеи.

    \textbf{Потребность пользователя}: получить готовое или песронализированное решение, которое превращает свободное время в осмысленное и эмоционально насыщенное "впечатление".

    \item[\large1.3] \large \textbf{Stakeholders (заинтересованные лица)}
    
    \normalsize
    \begin{itemize}[label=$-$]
        \item Конечные пользователи: потребители маршрутов и впечатлений
        \item Авторы впечатлений: экскурсоводы, эксперты, блогеры, энтузиасты
        \item Платформа: заинтересованные в развитии и монетизации продукта
        \item Администраторы проекта: отвечают за модерацию и качество впечатлений
        \item Партнёры: кафе, музеи, транспортные компании, потенциальные участиники маршрутов
    \end{itemize}

    \item[\large1.4] \large \textbf{Анализ конкурентов и существующих решений}
    \normalsize

    Существующие решения делятся на несколько категорий:

    \begin{enumerate}[label=\textbf{\arabic*.}]
        \item \textbf{Карты и навигационные приложения:} Google Maps, Яндекс Карты - решают задачу навигации, но не формируют "впечатления"
        \item \textbf{Туристические платформы:} TripAdvisor, Airbnb Experiences, GetYourGuide - ориентированы на туризм, часто перегружены и не адаптированы под спонтанный досуг
        \item \textbf{Социальные сети и блоги:} Instagram, YouTube, Телеграм каналы - предоставляют вдохновение, но не структурируют досуг, не персонализированны
    \end{enumerate}

    \textbf{Вывод:} на рынке отсутствует продукт, который объединяет маршруты, эмоции, персонализацию и реальное перемещение пользователя в одном сервисе.

    \item[\large1.5] \large \textbf{Вывод - существующие проблемы и актуальность разработки}
    
    \normalsize Существующие решения не отвечают на ключевой пользовательский вопрос: "что делать прямо сейчас и как провести время осмысленно?"

    Актуальность разработки обусловлена:
    \begin{itemize}[label=$-$]
        \item ростом спроса на короткие и спонтанные форматы досуга
        \item усталостью пользователей от типового туризма
        \item отсутствием платформы, ориентированной именно на "впечатления", а не точки на карте
    \end{itemize}

    \item[\Large2.] \Large \textbf{Скоуп проекта}
    
    \normalsize В продукт входит:

    \begin{itemize}[label=$-$]
        \item CRUD маршрутов
        \item CRUD впечатлений
        \item витрина впечатлений (какие есть впечатления, можно смотреть)
        \item покупка готовых маршрутов (можно и бесплатно, за рекламу например)
        \item realtime-отслеживание активного впечатления  (хочется подключить иишку для отслеживания местоположения)
        \item сервис рекомендаций (aka можно сделать что-то вроде скролла с рекомендациями впечатлений в зависимости от города, времени дня и тп)
        \item админка для управления контентом
        \item базовая аналитика пользовательской активности
    \end{itemize}

    В продукт не входит:

    \begin{itemize}[label=$-$]
        \item офлайн-навигация
        \item бронирование билетов и жилья (можно будет просто рекомендовать сторонние сервисы)
        \item социальная сеть и чаты между пользователями
    \end{itemize}

    Ограничения:

    \begin{itemize}[label=$-$]
        \item ограниченное количество городов из-за нехватки туристических данных
        \item отсутствие сложной ML-персонализации на старте
        \item минимальный набор авторов впечатлений на старте
        \item из-за новизны продкута - пользователь может испытывать предвзятость к использованию проекта
    \end{itemize}

    \item[\Large3.] \Large \textbf{Постановка цели проекта}

    \item[\large3.1] \large \textbf{Формулировка цели}

     \normalsize Цель проекта - создать цифровой продукт, который помогает пользователям быстро и удобно формировать персонализированный досуг в формате "впечатлений" во время поездок и прогулок, снижая когнитивную нагрузку при выборе маршрутов и активностей.

    \item[\large3.2] \large \textbf{Оценка достижения цели (метрики)}
    
    \normalsize
    \begin{itemize}[label=$-$]
        \item количество созданных и купленных впечатлений
        \item конверсия из просмотра витрины в покупку
        \item средняя длительность активного впечатления
        \item повторное использование приложения
        \item пользовательская оценка впечатлений
        \item реферальная программа и количество приглашённых друзей
    \end{itemize}

    \item[\large3.3] \large \textbf{Соответствие SMART}
    
    \normalsize
    \begin{itemize}[label=$-$]
        \item \textbf{Specific} - цель направлена на формирование досуга через впечатления
        \item \textbf{Measurable} - измеряется через конверсии, активность и retention
        \item \textbf{Achievable} - реализуема в рамках MVP
        \item \textbf{Relevant} - решает реальную пользовательскую проблему
        \item \textbf{Time-based} - достижение метрик оценивается в рамках этапов развития продукта
    \end{itemize}
\end{description}

\end{document}
